% SEGMENT OLP-0521-S01
% Part: counterfactuals
% Chapter: introduction
% Section: counterfactuals

\documentclass[../../../include/open-logic-section]{subfiles}

\begin{document}

\olfileid{cnt}{int}{cnt}

\olsection{உண்மைக்கு மாறான நிபந்தனைக் கூற்றுகள்}

ஆங்கிலத்தில் மிகவும் பொதுவானதும் முக்கியமானதுமான ``if \dots then
\dots'' கட்டமைப்புகளின் ஒரு வகை, \emph{to be} என்பதன் இறந்தகால
ஐயநிலை வடிவால் கட்டமைக்கப்படுகிறது: ``if it were the case that \dots
then it would be the case that \dots.'' இத்தகைய நிபந்தனையின் முன்னிலை
வழக்கமாகப் பொய்யாக, அதாவது உண்மைக்கு மாறாக இருப்பதால், இவை
\emph{உண்மைக்கு மாறான நிபந்தனைக் கூற்றுகள்} எனப்படுகின்றன. மேலும்,
\emph{to be} என்பதன் ஐயநிலை வடிவைப் பயன்படுத்துவதால் \emph{ஐயநிலை
நிபந்தனைக் கூற்றுகள்} என்றும் அழைக்கப்படுகின்றன. ``if it is the case
that \dots then it is the case that \dots'' என்ற வடிவமுள்ள
\emph{சுட்டுநிலை} நிபந்தனைக் கூற்றுகளிலிருந்து இவை வேறுபடுத்தப்படுகின்றன.
உண்மைக்கு மாறான நிபந்தனைகளும் சுட்டுநிலை நிபந்தனைகளும் வெவ்வேறு மெய்மை
நிபந்தனைகளைக் கொண்டுள்ளன. ஆடம்ஸின் புகழ்பெற்ற எடுத்துக்காட்டைக் கருதுக:
\begin{quote}
  ஆஸ்வால்ட் கென்னடியைக் கொல்லவில்லை எனில், வேறு யாரோ கொன்றார்.

  ஆஸ்வால்ட் கென்னடியைக் கொன்றிருக்காவிட்டால், வேறு யாரோ கொன்றிருப்பார்.
\end{quote}
முதலாவது சுட்டுநிலை நிபந்தனை; இரண்டாவது உண்மைக்கு மாறான நிபந்தனை.
முதலாவது தெளிவாக மெய்: அதிபர் ஜான் எஃப். கென்னடியை \emph{யாரோ}
கொன்றார் என்று நமக்குத் தெரியும். அந்த நபர் வாரன் அறிக்கைக்கு மாறாக லீ
ஹார்வி ஆஸ்வால்ட் அல்ல என்றால், வேறு யாரோ கென்னடியைக் கொன்றார். இரண்டாவது
வேறொன்றைக் கூறுகிறது. ஆஸ்வால்ட் கென்னடியைக் கொன்றிருக்காவிட்டால்---அதாவது
டாலஸில் நடந்த துப்பாக்கிச் சூடு தவிர்க்கப்பட்டோ தோல்வியடைந்தோ இருந்தால்---
பின்னர் வரலாறு வேறொரு படுகொலை வெற்றியடையும் விதத்தில் விரிந்திருக்கும்
என்று அது கூறுகிறது. அது மெய்யாக இருக்க, ஜேஎஃப்கே இறப்பதை உறுதிசெய்ய
வலிமைமிக்க சக்திகள் கூட்டுச்சதி செய்திருக்க வேண்டும்; ஜேஎஃப்கே படுகொலை
பற்றிய பல கூட்டுச்சதிக் கோட்பாட்டாளர்கள் இதையே நம்புகின்றனர்.

\emph{சுட்டுநிலை} நிபந்தனையைப் பொருள்சார் நிபந்தனை சரியாகப் பதிவு
செய்கிறதா என்பது இன்னும் விவாதத்திற்குரியது. குறிப்பாக, ஆங்கிலச்
சுட்டுநிலை நிபந்தனைகளின் மெய்மை நிபந்தனைகளைப் பொருள்சார் நிபந்தனை
தருகிறது என்பதுடன் பொருந்துமாறு, அதன் முரண்தோற்றங்களை ``விளக்க'' முடியுமா
என்பதும் விவாதத்திற்குரியது. இதற்கு மாறாக, உண்மைக்கு மாறான நிபந்தனைக்
கூற்றுகளைப் பொருள்சார் நிபந்தனையால் சரியாகக் குறியீடாக்க முடியாது என்பது
விவாதமற்றது. ஏனெனில், உண்மைக்கு மாறான நிபந்தனைகளின் முன்னிலைகள் பொதுவாகப்
பொய்யாக இருந்தாலும், பொய்யான முன்னிலையைக் கொண்ட அவை அனைத்தும் மெய்யல்ல.
எடுத்துக்காட்டாக, வாரன் அறிக்கையை நீங்கள் நம்பி, ஜேஎஃப்கே-யைக் கொல்லக்
கூட்டுச்சதி ஏதும் நடைபெறவில்லை என்றும் நம்பினால், ஆடம்ஸின் உண்மைக்கு
மாறான நிபந்தனையே ஓர் எதிரெடுத்துக்காட்டு.

காரணவியல் பகுத்தறிதலில் உண்மைக்கு மாறான நிபந்தனைக் கூற்றுகள் முக்கியப்
பங்கு வகிக்கின்றன: காரணவியல் தொடர்புகளை வெளிப்படுத்துவது அவற்றின் ஒரு
முதன்மைப் பயன்பாடு. எடுத்துக்காட்டாக, ஒரு தீக்குச்சியை உரசுவது அது
பற்றுவதற்குக் காரணமாகிறது; ``இந்தத் தீக்குச்சி உரசப்பட்டிருந்தால், அது
பற்றியிருக்கும்'' என்று இதை வெளிப்படுத்தலாம். பொருள்சார் நிபந்தனைகளாலும்,
பொதுவாகச் சுட்டுநிலை நிபந்தனைகளாலும் இதை வெளிப்படுத்த முடியாது:
``தீக்குச்சி உரசப்படுகிறது $\lif$ தீக்குச்சி பற்றுகிறது'' என்பது,
தீக்குச்சி ஒருபோதும் உரசப்படாவிட்டால், அது உரசப்பட்டிருந்தால் என்ன
நடந்திருக்கும் என்பதைப் பொருட்படுத்தாமல் மெய்யாகும். இதைவிட மோசமாக,
``தீக்குச்சி உரசப்படுகிறது $\lif$ தீக்குச்சி மலர்க்கொத்தாக மாறுகிறது''
என்பதும் தீக்குச்சி ஒருபோதும் உரசப்படாவிட்டால் மெய்யாகும்; ஆனால் அது
உரசப்பட்டிருந்தால் உறுதியாக மலர்க்கொத்தாக மாறியிருக்காது.

உண்மைக்கு மாறான நிபந்தனைகளின் சரியான தருக்கம் துல்லியமாக எது என்பது
இன்னும் விவாதிக்கப்படுகிறது. ஸ்டால்நேக்கரும் லூயிஸும் அவற்றைப் பற்றிய
செல்வாக்குமிக்க பகுப்பாய்வு ஒன்றை வழங்கினர். அவர்களின்படி, ``$S$ ஆக
இருந்திருந்தால் $T$ ஆக இருந்திருக்கும்'' என்ற உண்மைக்கு மாறான நிபந்தனை,
நடப்பு உலகம் அமைந்துள்ள விதத்துக்கு மிக நெருக்கமாகவும் $S$ மெய்யாகவும்
இருக்கும் உண்மைக்கு மாறான சூழ்நிலையில் (``சாத்தியமான உலகில்'') $T$
மெய்யானால், எனில் மட்டுமே, மெய்யாகும். சாத்தியமான உலகங்களின் இருப்பியலை
இது மேற்கோள்கொள்வதால், இது ``இருப்பியல்'' பகுப்பாய்வு எனப்படுகிறது.
மற்ற பகுப்பாய்வுகள் நிபந்தனை நிகழ்தகவுகளையோ நம்பிக்கைத் திருத்தக்
கோட்பாடுகளையோ பயன்படுத்துகின்றன. உண்மைக்கு மாறான நிபந்தனைகளுக்கு
முன்மொழியப்பட்ட பல்வேறு தருக்கங்கள் பெருகியுள்ளன. லூயிஸ்--ஸ்டால்நேக்கர்
தருக்கம் என்று ஒன்றே ஒன்றுகூட இல்லை: மிக நெருங்கிய சாத்தியமான உலகங்களை
மேற்கோள்கொண்டு ஸ்டால்நேக்கரும் லூயிஸும் ஒத்த வழிகளில் விளக்கங்களை
முன்மொழிந்தாலும், அவர்கள் ஏற்ற எடுகோள்கள் வெவ்வேறு செல்லுபடியான
அனுமானங்களைத் தருகின்றன.

\end{document}
